\section{Experiments}
\label{sec:experiments}

\subsection{Experimental setup}

We assemble ten image-to-video systems: HunyuanVideo 1.5, HunyuanVideo 1.5 Distill, MiniMax H3, Kling 3.0 Standard, Seedance 2.5, Veo 3.1 Fast, CogVideoX 1.5, Wan 2.1, Wan 2.2, and Wan 3.0. Each system has the same 500 reference-image/task pairs and frozen generation-facing prompts, yielding 5,000 published videos. Videos are mapped by task identifier and decoded under a common evaluation pipeline. A previously evaluated candidate set is excluded from publication because its generation provenance cannot be verified independently; none of its videos or scores contribute to the results below.

The formal judge is Qwen3-VL-235B-A22B-Instruct-FP8, served with four-GPU tensor parallelism. Four deterministic data shards are evaluated concurrently. Every released model collection contains exactly 500 videos; Seedance task \texttt{erob\_190} is now present. On September 11, 2026, 451 videos across nine collections were refreshed (50 per collection and 51 for Seedance). Under the frozen-cache policy, this invalidates pre-refresh scores for the affected collections and requires complete re-evaluation before publication.

\subsection{Main results}

Table~\ref{tab:main-results} reports results confirmed by canonical aggregate JSON files. A dash denotes a system without a complete aggregate; partial-sample scores are not substituted for complete results.

\begin{table*}[t]
\tablestyle{4pt}{1.1}
\caption{Current evaluation status on the refreshed 500-task operational split. Dashes indicate that post-refresh complete aggregates are pending; superseded pre-refresh values are intentionally omitted.}
\label{tab:main-results}
\resizebox{\textwidth}{!}{%
\begin{tabular}{lrrrrrrr}
\toprule
Model & \shortstack{Industrial logic and\\fact alignment} & \shortstack{Geometric\\integrity} & \shortstack{Physical\\plausibility} & \shortstack{Temporal\\consistency} & \shortstack{Reference and\\motion fidelity} & \shortstack{Application\\usefulness} & \shortstack{Ranking\\score} \\
\midrule
HunyuanVideo 1.5 & -- & -- & -- & -- & -- & -- & -- \\
HunyuanVideo 1.5 Distill & -- & -- & -- & -- & -- & -- & -- \\
MiniMax H3 & -- & -- & -- & -- & -- & -- & -- \\
Kling 3.0 Standard & -- & -- & -- & -- & -- & -- & -- \\
Seedance 2.5 & -- & -- & -- & -- & -- & -- & -- \\
Veo 3.1 Fast & -- & -- & -- & -- & -- & -- & -- \\
CogVideoX 1.5 & -- & -- & -- & -- & -- & -- & -- \\
Wan 2.1 & -- & -- & -- & -- & -- & -- & -- \\
Wan 2.2 & -- & -- & -- & -- & -- & -- & -- \\
Wan 3.0 & -- & -- & -- & -- & -- & -- & -- \\
\bottomrule
\end{tabular}
}
\end{table*}

No post-refresh model currently has a complete aggregate that is valid for publication. The table therefore reports no point estimates or confidence intervals. Historical aggregates produced before the video refresh remain useful only as pipeline diagnostics and are not comparable to the current Hugging Face release. Rankings, cluster-bootstrap confidence intervals, and axis-level comparisons will be restored only after every model has been evaluated against the current 500-file inventory and passes \texttt{ranking\_publishable=true}.

\paragraph{Analysis framework.}
The analysis asks whether generic visual coherence guarantees industrial performance, how scores change with difficulty, and which capability-specific trade-offs are hidden by a single mean.

\subsection{Domain and task breakdown}

Domain--task analysis uses a $5\times5$ matrix and macro averages over populated cells. It distinguishes general geometry weaknesses from failures concentrated in periodic precision structures and separates fluid-evolution errors from causal incompleteness in emergency scenes.

\subsection{Difficulty and application breakdown}

Results are stratified by the four preassigned difficulty tiers and six application types. The analysis considers not only score decline but also which axis first becomes limiting as task complexity rises. Application reporting separately measures required-event coverage and strict application pass rate.

\subsection{Qualitative failure analysis}

Qualitative analysis focuses on cases where conventional visual quality and \FORGE{} disagree: (i) smooth but static clips that omit required camera motion; (ii) visually dramatic failures with impossible topology; (iii) plausible single frames with broken temporal identity; (iv) localized defect prompts that trigger global regeneration; and (v) safety scenarios missing a trigger, response, or terminal state. Reference images, chronological evidence frames, per-axis scores, gate reasons, and judge evidence jointly localize these industrial failures.

\subsection{Limitations}

The benchmark is finite and does not cover every industrial sector, geographic standard, or equipment type. Some physical properties cannot be inferred reliably from pixels alone. Multimodal judges can be biased or inconsistent even under deterministic decoding, and contact-sheet evidence compresses temporal information. The operational split is balanced by domain rather than by real deployment frequency. Finally, benchmark scores indicate conformance to annotated visual contracts, not certification of safety, engineering correctness, or legal compliance.
